\documentclass[a4paper]{article}
\usepackage{a4wide}
\usepackage{amsmath}

\title{Ideas for a High-Level FIRRTL semantics in Coq}
\author{David N. Jansen}

\begin{document}
\maketitle

This note builds upon the file \texttt{Firrtl.v} written by Xiaomu Shi
and proposes a plan to extend the semantics to high-level FIRRTL
in a way that catches the ``last connect wins'' correctly.
(In low-level FIRRTL, every signal is allowed to be connected only once,
so this principle is not necessary.)

An internal module in FIRRTL consists of a name, a list of ports, and a single statement.
This statement typically is a statement group.

The translation of a module first reads all ports (input and output) and defines based on that the type $\mathtt{input}$ and $\mathtt{output}$ of the translation of the module.

The translation of the statements of a module is a tuple consisting of
\begin{description}

\item[The state,]
	a description of the registers and memory defined in the module, together with their types.
	The state can probably be modelled as the \emph{type of} a value store.

\item[The wires] defined in the module.
	They can be modelled similar to states, with the difference that wires are not stored over time.

\item[The nodes] defined in the module.
	A node is just a shorthand for a certain expression;
	likely this is a value store.

\item[The transition function]
	calculates the new value assigned to registers/memories upon the next rising clock edge, together with the output signals.
	The type of the transition function is something like
	\texttt{tr :} \textit{readables} $\to$ \textit{writables}.

	The readables are the input signals, wires, nodes, registers and memories;
	the writables are wires, output signals, and new values of registers and memories.
\end{description}
A single statement transforms this tuple to another one.
\begin{description}

\item[The empty statement] \texttt{skip} applies identity transform.

\item[A wire definition] \texttt{wire} \textit{name}\texttt{:} \textit{type}
	extends the wires by $\mathit{name}$.

\item[A node definition]
	extends the nodes, similar to a wire definition.

\item[A register definition] $\mathtt{reg~}\mathit{name}\mathtt{: }\mathit{type}$
	extends $\mathtt{state}$ by $\mathit{name}$
	and initializes $(\mathtt{newtr~}\mathit{input}\mathtt{~}\mathit{state})_\mathit{name}$ to $\mathit{name}$.

	(Page 24, subsection 5.9.3 of the FIRRTL specification:
	``Registers do not need to be connected to under all conditions,
	as it will keep its previous value if unconnected.'')

\item[A memory definition] has a similar effect as register definitions.

\item[A full connect] $\mathit{var}\mathtt{~<=~}\mathit{expression}$
	replaces function $\mathtt{tr}$ by the one where $\mathit{var}$ is assigned $\mathit{expression}$ and for the rest is identical with $\mathtt{tr}$.

\item[A conditional statement]
	\texttt{when} \textit{condition} \texttt{:} \textit{true-statement} \texttt{else :} \textit{false-statement}.
	Note that components (in particular registers and memories, but also wires and nodes) defined in the branches
	should always be included in the result.
	Connections to such objects are always made, even when the condition is not satisfied.

	To translate this statement, first find $\mathtt{state}_\mathit{true}$ and $\mathtt{tr}_\mathit{true}$, the transformation of $\mathtt{state}$ and $\mathtt{tr}$ under $\textit{\textmd{true-statement}}$,
	and $\mathtt{state}_\mathit{true}$ and $\mathtt{tr}_\mathit{false}$, the corresponding transformation under $\textit{\textmd{false-statement}}$.
	Then we set
	\begin{align*}
	\mathtt{newstate} & := \mathtt{state}_\mathit{true} \cup \mathtt{state}_\mathit{false} \\
	\mathtt{newtr~}\mathit{input}\mathtt{~}\mathit{state}\mathtt{~}v & :=
	\begin{cases}
		\mathtt{tr}_\mathit{true} \mathtt{~}v & \text{if } v \in \mathtt{state}_\mathit{true} \setminus \mathtt{state} \\
		\mathtt{tr}_\mathit{false}\mathtt{~}v & \text{if } v \in \mathtt{state}_\mathit{false} \setminus \mathtt{state} \\
		\mathtt{mux~}\mathit{condition}\mathtt{~(tr}_\mathit{true}\mathtt{~}v\mathtt{)~(tr}_\mathit{false}\mathtt{~}v\mathtt{)} & \text{if } v \in \mathtt{state}
	\end{cases}
	\end{align*}
	Wires and nodes defined in the branches are handled similarly.
	I think it is an error if both branches define the same element (or different elements with the same name);
	we should have $\mathtt{state}_\mathit{true} \cap \mathtt{state}_\mathit{false} = \mathtt{state}$.

\item[Statement groups]
	are translated by transforming the state, named signals and transition function in order.
	Because a later connect overwrites an earlier one,
	we automatically get ``last connect wins''.
\end{description}

To translate the statements in a module,
one initializes the state, wires, and nodes to empty, and the transition function to the function that is undefined everywhere,
and then transforms these by the statement (group) in the module.
\end{document}
